\begin{proposition}[p186の系7.1.6]
	$A \subseteq \nat$とする.\\
  1. $A$は可算集合である. 
  さらに,
  次の条件(1)と(2)は同値である.\\
  (1)$A \subseteq n$をみたす自然数$n$が存在する.\\
  (2)$A$は有限集合である.\\
  論理式は以下.
  \begin{equation*}
     \exists n (A \subseteq n) \leftrightarrow
     \exists n \in \nat \exists f (\BijectionMap{f}{n}{A})
  \end{equation*}
  2. $A$が空でなければ$\seigen{C_A}{A}$による,
  $0 \in C_A[A]$の逆像$l \in A$は$A$の最小限である.
  論理式は以下.
  \begin{equation*}
     A \neq \emptyset \rightarrow
     \exists l (\seigen{C_A}{A}(l) = 0 \land \forall n \in A(l \leq n))
  \end{equation*}
\end{proposition}

